\documentclass[aps,pra,twocolumn,superscriptaddress,longbibliography,floatfix]{revtex4-2}

\usepackage{amsmath,amssymb,amsthm,graphicx,xcolor,bm,braket,dcolumn,booktabs}
\usepackage{hyperref}

\newtheorem{theorem}{Theorem}

%% ========================================================================
\title{Composed Natural Geometries for Core-Valence Diatomics:\\
LiH and BeH$^+$ from Discrete Graph Topology}
\thanks{Paper 17 in the Geometric Vacuum series}

\author{J.~Loutey}
\affiliation{Independent Researcher, Kent, Washington}

\date{March 21, 2026}

%% ========================================================================
%% ABSTRACT
%% ========================================================================
\begin{document}
\begin{abstract}
The natural geometry hierarchy assigns each separable quantum system
a coordinate system where Coulomb singularities become boundary
conditions:\ $S^3$ for atoms (Paper~7), prolate spheroids for
H$_2^+$ (Paper~11), hyperspherical for He (Paper~13), molecule-frame
hyperspherical for H$_2$ (Paper~15).
Systems with core electrons resist a single natural geometry because
core and valence electrons occupy different length scales.
We introduce \emph{composed natural geometries}---fiber bundles where
each electron group inhabits its own natural coordinate system,
connected by screening functions and ab initio Phillips--Kleinman
pseudopotentials derived entirely from the core wavefunction.
For LiH ($Z_A=3$, $Z_B=1$, 4~electrons): the Li$^+$~$1s^2$ core is
solved via Level~3 hyperspherical (Paper~13), yielding $Z_{\rm eff}(r)$
and pseudopotential parameters; the two valence electrons are solved via
Level~4 molecule-frame hyperspherical (Paper~15) with $Z_{\rm eff}$
injection; the potential energy surface is scanned via a $\rho$-collapse
adiabatic method that reduces per-$R$-point cost to seconds.
Results: $R_{\rm eq} = 3.21$~bohr (6.4\% error vs.\ experiment 3.015,
improved from 17\% in LCAO),
$\omega_e = 1471$~cm$^{-1}$ (4.6\% error vs.\ 1406),
with zero experimental molecular data in the pseudopotential.
Total wall time is ${\sim}4$~minutes per molecule on commodity hardware.
The same formula transfers to BeH$^+$ ($Z_A=4$) with no per-molecule
tuning.
\end{abstract}

\maketitle

%% ========================================================================
%% I. INTRODUCTION
%% ========================================================================
\section{Introduction}
\label{sec:intro}

The Geometric Vacuum program constructs discrete lattice
Hamiltonians by identifying the \emph{natural geometry} of each
quantum system and discretizing it as a sparse graph~\cite{paper0,paper1,paper7}.
Four levels of this hierarchy have been established:

\begin{description}
\item[Level~1] \emph{One-center, one-electron} (H).
  Fock's stereographic projection~\cite{fock1935} maps the hydrogen atom
  to the unit $S^3$, where the $1/r$ potential becomes the conformal
  factor.  The discrete graph Laplacian reproduces hydrogenic energies
  to $< 0.1\%$ (Paper~7~\cite{paper7}).

\item[Level~2] \emph{Two-center, one-electron} (H$_2^+$).
  Prolate spheroidal coordinates separate the two-center Coulomb problem
  exactly.  Paper~11~\cite{paper11} achieves 0.70\% error with a sparse
  lattice and zero free parameters.

\item[Level~3] \emph{One-center, two-electron} (He).
  Hyperspherical coordinates $(R_e, \alpha, \hat{r}_1, \hat{r}_2)$
  resolve the $e$--$e$ cusp as a boundary condition at $\alpha=\pi/4$,
  $\theta_{12}=0$.  Paper~13~\cite{paper13} achieves 0.05\% error for
  He with a single adiabatic channel.

\item[Level~4] \emph{Two-center, two-electron} (H$_2$).
  Molecule-frame hyperspherical coordinates
  $(R_e, \alpha, \theta_1, \theta_2, \Phi)$ parameterized by
  $\rho = R/(2R_e)$ resolve all five Coulomb singularities simultaneously.
  Paper~15~\cite{paper15} recovers 94.1\% of $D_e$ for H$_2$ and
  93.1\% for HeH$^+$ with $\sigma{+}\pi$ channels.
\end{description}

Level~4 handles H$_2$ and HeH$^+$ but these are \emph{core-free}
systems: all electrons participate in bonding.  LiH has 2~core
electrons on Li and 2~valence electrons, occupying different length
scales.  The LCAO approach~\cite{fci_m} gives $R_{\rm eq} = 2.5$~bohr
(17\% error vs.\ experiment 3.015) because the graph Laplacian's
intra-atom kinetic energy is $R$-independent---the atomic lattice
structure does not change as atoms approach~\cite{fci_m}.
Paper~15 identified this as the ``locked-core boundary'' problem:
a single Level~4 geometry for all four electrons would require $S^{11}$
with ${\sim}10^6$ angular basis functions, computationally infeasible.

We solve this by \emph{composing} geometries from different levels.
Each electron group gets its own natural coordinate system.
The inter-group coupling is handled by three typed edges:
(i)~screening via $Z_{\rm eff}(r)$ from the core density,
(ii)~Pauli exclusion via a Phillips--Kleinman~\cite{pk1959}
pseudopotential, and
(iii)~nuclear repulsion via bare charges.
This is the first realization of a \emph{category} of natural
geometries with morphisms preserving lattice structure.

The key results: LiH $R_{\rm eq} = 3.21$~bohr (6.4\% error),
$\omega_e = 1471$~cm$^{-1}$ (4.6\% error), with zero experimental
molecular data.  The same code, with only $(Z_A, Z_B, M_A, M_B)$
changed, produces a bound BeH$^+$ with physical spectroscopic
constants.  We do not claim production accuracy; this is a proof of
concept for the composed geometry architecture.


%% ========================================================================
%% II. THEORY
%% ========================================================================
\section{Theory: Composed Natural Geometries}
\label{sec:theory}

\subsection{The composition principle}
\label{sec:composition}

The composed graph for a core-valence diatomic is
\begin{equation}
  G_{\rm total} = G_{\rm nuc} \ltimes G_{\rm core}(R)
    \ltimes G_{\rm val}(R, \text{core state})\,,
  \label{eq:Gtotal}
\end{equation}
where $\ltimes$ denotes the semi-direct product (fiber bundle).
Each layer inhabits a different natural geometry:

\begin{description}
\item[Layer~0 (nuclear):]
  Morse $\text{SU}(2) \times \text{SO}(3)$ from Paper~10~\cite{paper10}.
  States $\ket{v, J, M_J}$.

\item[Layer~1 (core):]
  Level~3 hyperspherical on the heavy atom at charge $Z_A$.
  States parameterized by $(R_e^{\rm core}, \alpha_{\rm core})$.

\item[Layer~2 (valence):]
  Level~4 molecule-frame hyperspherical.
  States parameterized by
  $(R_e^{\rm val}, \alpha, \theta_1, \theta_2, \Phi; \rho)$.
\end{description}

Backward compatibility between levels fails for a fundamental reason:
Level~4 coordinates require two electrons (hyperradius
$R_e = \sqrt{r_1^2 + r_2^2}$, hyperangle $\alpha = \arctan(r_2/r_1)$),
which are undefined for a single electron.  Level~4 cannot subsume
Level~1.  Similarly, Level~3 requires two electrons and cannot
represent one.  The hierarchy is a two-dimensional grid
(centers~$\times$~electrons), not a linear tower.  Composition is
necessary because no single coordinate system handles all electron
groups simultaneously.

The total energy at internuclear distance $R$ is
\begin{equation}
  E_{\rm total}(R) = E_{\rm core} + V_{\rm cross}(R)
    + E_{\rm elec}(R) + V_{NN}(R)\,,
  \label{eq:Etotal}
\end{equation}
where $E_{\rm core}$ is the $R$-independent core energy from the
Level~3 solve, $V_{\rm cross}(R)$ is the core--$B$ nuclear attraction,
$E_{\rm elec}(R)$ is the valence electronic energy from the Level~4
solve, and $V_{NN}(R) = Z_A Z_B / R$ is the bare nuclear repulsion.

The core-to-other-nucleus attraction for a $1s^2$ core is computed
analytically:
\begin{equation}
  V_{\rm cross}(R) = -n_{\rm core}\, Z_B\,
    \Braket{1s_{Z_A} | \frac{1}{r_B} | 1s_{Z_A}}\,,
  \label{eq:Vcross}
\end{equation}
where the expectation value evaluates to
\begin{equation}
  \Braket{1s_Z | \frac{1}{r_B} | 1s_Z}
    = \frac{1}{R}\bigl[1 - (1 + ZR)\,e^{-2ZR}\bigr]\,.
  \label{eq:1s_expectation}
\end{equation}
This uses a hydrogenic $1s$ approximation for the core orbital; a more
accurate treatment would integrate $Z_B/r_B$ against the numerically
computed core density from the Level~3 solve.  For Li$^+$ ($Z=3$),
the hydrogenic approximation is accurate to ${\sim}1\%$ because the
two-electron core is well-described by an effective nuclear charge
$Z_{\rm eff}^{\rm core} \approx 2.7$.


\subsection{Core screening function}
\label{sec:screening}

The core is solved via the Level~3 adiabatic hyperspherical
method~\cite{lin1995} (Paper~13~\cite{paper13}) at nuclear charge $Z_A$.  The one-electron
density $n(r)$ is extracted by integrating $|\psi(R_e, \alpha)|^2$
over the hyperangle and the coordinates of the second electron,
using $r_1 = R_e \cos\alpha$, $r_2 = R_e \sin\alpha$.
At each hyperradius $R_e$, the angular eigenfunction is resolved on a
finite-difference grid of $n_\alpha = 200$ points, and contributions
from both electrons are binned onto a radial grid.

The cumulative electron count
\begin{equation}
  N_{\rm core}(r) = \int_0^r n(r')\,dr'
  \label{eq:Ncore}
\end{equation}
gives the effective nuclear charge
\begin{equation}
  Z_{\rm eff}(r) = Z_A - N_{\rm core}(r)\,,
  \label{eq:Zeff}
\end{equation}
with boundary conditions $Z_{\rm eff}(r\to 0) = Z_A$ (no screening at
the nucleus) and $Z_{\rm eff}(r\to\infty) = Z_A - 2$ (full screening
by two electrons).  For Li$^+$: $Z_{\rm eff}(0) = 3.0$,
$Z_{\rm eff}(2\text{ bohr}) \approx 1.01$.


\subsection{Ab initio Phillips--Kleinman pseudopotential}
\label{sec:pk}

The valence electrons must be orthogonal to the core.  In continuous
quantum chemistry, this is enforced by the
Phillips--Kleinman~\cite{pk1959} pseudopotential.
We parameterize the PK barrier as
\begin{equation}
  V_{\rm PK}(r) = \frac{A\,e^{-Br^2}}{r^2}\,,
  \label{eq:Vpk}
\end{equation}
where $A$ and $B$ are derived entirely from the core wavefunction.

\paragraph{Barrier width $B$.}
The PK barrier must be localized where the core density is concentrated.
We define the inverse-square-weighted core radius
\begin{equation}
  r_{\rm eff} = \sqrt{\frac{\int n(r)\,dr}
    {\int n(r)/r^2\,dr}}\,,
  \label{eq:reff}
\end{equation}
which weights the inner core---where $1/r^2$ is large---more heavily
than the diffuse tail.  Then
\begin{equation}
  B = \frac{1}{2\,r_{\rm eff}^2}\,.
  \label{eq:B}
\end{equation}
For Li: $r_{\rm eff} = 0.267$~bohr, $B = 7.00$~bohr$^{-2}$.
For Be: $r_{\rm eff} = 0.200$~bohr, $B = 12.53$~bohr$^{-2}$.

\paragraph{Barrier height $A$.}
The kinetic energy cost of core--valence orthogonalization is
\begin{equation}
  A = \bigl|E_{\rm core}/N_{\rm core} - E_{\rm val}\bigr|
      \cdot N_{\rm core}\,,
  \label{eq:A}
\end{equation}
where $E_{\rm core}$ is the total core energy from the hyperspherical
solve, $N_{\rm core} = 2$, and $E_{\rm val}$ is estimated from the
first ionization energy of the neutral atom (an atomic property, not
a molecular one: $\text{IE}(\text{Li}) = 0.198$~Ha~\cite{nist},
$\text{IE}(\text{Be}) = 0.343$~Ha~\cite{nist}).
For Li: $A = 6.93$~Ha$\cdot$bohr$^2$.
For Be: $A = 13.01$~Ha$\cdot$bohr$^2$.

No experimental molecular data enters these formulas.  Only atomic
properties: $Z$, the core wavefunction (via the Level~3 solve), and
the first ionization energy of the neutral atom.


\subsection{\texorpdfstring{$Z_{\rm eff}$}{Z\_eff} injection into the
Level~4 solver}
\label{sec:injection}

The Level~4 angular Hamiltonian (Paper~15~\cite{paper15}) contains
nuclear attraction terms $-Z_A/r_{1A} - Z_A/r_{2A}$.  We replace
the constant $Z_A$ with the function $Z_{\rm eff}(r)$ from the core
solution.  In the coupled-channel expansion, the nuclear coupling
matrix element between channels $(l_1', l_2')$ and $(l_1, l_2)$ at
fixed $\alpha$ involves a Gaunt integral weighted by the nuclear
potential.  For constant $Z_A$, these are computed via an analytical
multipole expansion.  For $Z_{\rm eff}(r)$, we add a smooth
Gauss--Legendre quadrature correction for the
$[Z_A - Z_{\rm eff}(r)]/r$ difference.

Since $Z_{\rm eff}(0) = Z_A$, the correction vanishes at $r \to 0$,
eliminating the $1/r$ singularity from the quadrature.  This ensures
machine-precision backward compatibility when $Z_{\rm eff}$ is
constant.


\subsection{The \texorpdfstring{$\rho$}{rho}-collapse and fast PES
scanning}
\label{sec:rho_collapse}

Paper~15 proved that the Level~4 angular eigenvalue problem depends
on $(R, R_e)$ only through the dimensionless ratio
$\rho = R/(2R_e)$.  This means the angular eigenvalue curves
$\varepsilon_\nu(\rho)$ can be computed once on a one-dimensional grid
of $\rho$ values and then interpolated via cubic splines.

The effective radial potential is
\begin{equation}
  U(R_e; R) = \frac{\varepsilon_0(\rho) + 15/8}{R_e^2}\,,
  \quad \rho = \frac{R}{2R_e}\,,
  \label{eq:U_eff}
\end{equation}
and the electronic energy at each $R$ is obtained by solving the
one-dimensional radial Schr\"odinger equation
\begin{equation}
  \Bigl[-\tfrac{1}{2}\frac{d^2}{dR_e^2} + U(R_e; R)\Bigr]
    F(R_e) = E_{\rm elec}\,F(R_e)
  \label{eq:radial}
\end{equation}
via a tridiagonal eigensolver (cost: milliseconds).

The angular cache is built on a non-uniform $\rho$ grid of
$n_\rho \approx 80$ points (denser at small $\rho$ where the
eigenvalue changes rapidly), with the static Hamiltonian (kinetic,
centrifugal, Liouville, $e$--$e$ coupling) precomputed once and the
nuclear coupling added at each $\rho$.  Cache construction for a
single $R$-point takes ${\sim}2$~seconds.


%% ========================================================================
%% III. COMPUTATIONAL DETAILS
%% ========================================================================
\section{Computational Details}
\label{sec:details}

All calculations use the following parameters unless stated otherwise.

\paragraph{Core solve.}
The Li$^+$ $1s^2$ core (or Be$^{2+}$ $1s^2$) is solved via the
Level~3 adiabatic hyperspherical method with
$l_{\max} = 2$, $n_\alpha = 200$, $N_{R}^{\rm angular} = 200$
$R$-points for the adiabatic curves, and
$N_{R}^{\rm radial} = 2000$ grid points for the radial solve.
Wall time: ${\sim}3$~seconds.

\paragraph{Angular cache.}
$n_\rho = 80$, $\rho \in [0.02, 5.0]$, $l_{\max} = 2$
($\sigma$ channels only; $n_{\rm ch}$ determined by $l_{\max}$).
$n_\alpha = 100$ grid points in the angular finite-difference mesh.
Cache build per $R$-point: ${\sim}2$~seconds.

\paragraph{PES scan.}
18 $R$-points, $R \in [2.0, 7.0]$~bohr, with denser sampling near
the minimum.  Per-$R$ radial solve: $n_{R_e} = 300$ grid points,
$R_e \in [0.3, 15.0]$~bohr, tridiagonal eigensolver.
Total PES scan time: ${\sim}48$~seconds.

\paragraph{Morse fit.}
Nonlinear least-squares fit (Levenberg--Marquardt) to the Morse
potential $E(R) = E_{\min} + D_e[1 - e^{-a(R-R_{\rm eq})}]^2$
within $\pm 1.5$~bohr of the grid minimum.

\paragraph{Reduced masses.}
$\mu(\text{LiH}) = 1822.888 \times 7.016 \times 1.008 / (7.016 + 1.008)$~a.u.;
$\mu(\text{BeH}^+) = 1822.888 \times 9.012 \times 1.008 / (9.012 + 1.008)$~a.u.

\paragraph{Total pipeline time.}
${\sim}4$~minutes per molecule on a single-core commodity workstation.

\paragraph{What is not included.}
Core--valence exchange integrals (deferred),
non-adiabatic coupling between angular channels (DBOC only),
$l_{\max} > 2$ (results at $l_{\max} = 3$ cited for comparison),
$\pi$ channels ($\sigma$-only at $l_{\max} = 2$).

All calculations were performed with GeoVac
v1.6.1~\cite{geovac2026}.  The complete pipeline (core solve, angular
cache, PES scan, Morse fit, nuclear lattice) is implemented in the
class \texttt{ComposedDiatomicSolver} and is reproducible with the
command
\texttt{ComposedDiatomicSolver.LiH\_ab\_initio(l\_max=2).run\_all()}.


%% ========================================================================
%% IV. RESULTS
%% ========================================================================
\section{Results}
\label{sec:results}

\subsection{Ab initio pseudopotential parameters}
\label{sec:pk_params}

Table~\ref{tab:pk} collects the PK parameters for Li and Be cores,
derived entirely from the core wavefunction and atomic ionization
energies.

\begin{table}[b]
\caption{Ab initio Phillips--Kleinman pseudopotential parameters.
  All quantities from the core wavefunction; no molecular data.}
\label{tab:pk}
\begin{ruledtabular}
\begin{tabular}{lddd}
\textrm{Quantity} & \multicolumn{1}{c}{Li ($Z{=}3$)}
  & \multicolumn{1}{c}{Be ($Z{=}4$)}
  & \multicolumn{1}{c}{Formula} \\
\colrule
$r_{\rm eff}$ (bohr)   & 0.267  & 0.200  & \multicolumn{1}{c}{Eq.~(\ref{eq:reff})} \\
$A$ (Ha$\cdot$bohr$^2$) & 6.93  & 13.01  & \multicolumn{1}{c}{Eq.~(\ref{eq:A})} \\
$B$ (bohr$^{-2}$)       & 7.00  & 12.53  & \multicolumn{1}{c}{Eq.~(\ref{eq:B})} \\
$E_{\rm core}$ (Ha)     & -7.280 & -13.700 & \multicolumn{1}{c}{Level~3} \\
\end{tabular}
\end{ruledtabular}
\end{table}

$B$ scales approximately as $Z^2$ ($7.00/12.53 \approx 9/16 = (3/4)^2$),
consistent with the core radius scaling as $1/Z$.


\subsection{LiH potential energy surface}
\label{sec:lih}

Table~\ref{tab:lih} compares spectroscopic constants for LiH under
different pseudopotential treatments.

\begin{table}
\caption{LiH spectroscopic constants at $l_{\max}=2$.
  ``LCAO'' is the atom-centered Full CI result from Ref.~\cite{fci_m}.
  ``No PK'' disables the pseudopotential; ``Ab initio PK'' uses
  Eqs.~(\ref{eq:A})--(\ref{eq:B}); ``Fitted PK'' uses manually
  chosen $A=5.0$, $B=7.0$.  Experimental data from
  Ref.~\cite{huber1979}.}
\label{tab:lih}
\begin{ruledtabular}
\begin{tabular}{lddddd}
\textrm{Property}
  & \multicolumn{1}{c}{LCAO}
  & \multicolumn{1}{c}{No PK}
  & \multicolumn{1}{c}{Ab initio}
  & \multicolumn{1}{c}{Fitted}
  & \multicolumn{1}{c}{Expt.} \\
\colrule
$R_{\rm eq}$ (bohr)      & 2.5   & 0.88  & 3.21  & 2.97  & 3.015 \\
$\omega_e$ (cm$^{-1}$)   & \multicolumn{1}{c}{---}  & 2637  & 1471  & 1617  & 1406 \\
$B_e$ (cm$^{-1}$)        & \multicolumn{1}{c}{---}  & \multicolumn{1}{c}{---}  & 6.63  & 7.84  & 7.51 \\
Bound?                    & \multicolumn{1}{c}{Yes}
  & \multicolumn{1}{c}{collapse}
  & \multicolumn{1}{c}{Yes}
  & \multicolumn{1}{c}{Yes}
  & \multicolumn{1}{c}{Yes} \\
\end{tabular}
\end{ruledtabular}
\end{table}

Without the PK pseudopotential, the valence electrons collapse into
the core ($R_{\rm eq} = 0.88$~bohr).  The PK barrier is necessary, not
merely helpful.  The ab initio PK gives $\omega_e$ closer to experiment
(4.6\% error) than the fitted PK (15\% error), suggesting that the
$\langle 1/r^2 \rangle$-weighted barrier width captures the correct
PES curvature.
The ab initio $R_{\rm eq}$ overshoots experiment (3.21 vs.\ 3.015~bohr)
because the ab initio barrier height $A = 6.93$ exceeds the fitted
value $A = 5.0$ by 39\%, producing stronger Pauli repulsion that
pushes the equilibrium outward.  The barrier widths are nearly
identical ($B = 7.00$ ab initio vs.\ $7.0$ fitted), so the overshoot
is controlled entirely by $A$.  Reducing $A$ toward 5 would recover
the fitted $R_{\rm eq}$ at the cost of reintroducing a fitted
parameter; we retain the ab initio value as the parameter-free
prediction.

The $l_{\max}$ dependence provides a convergence indicator:
$l_{\max} = 2$ gives $R_{\rm eq} = 2.97$~bohr (fitted) or 3.21~bohr
(ab initio); $l_{\max} = 3$ gives $R_{\rm eq} = 3.23$~bohr (fitted).
The experimental value (3.015) is bracketed by $l_{\max} = 2$ and
$l_{\max} = 3$ with fitted PK, and within 6.4\% at $l_{\max} = 2$
with the ab initio PK.

$D_e$ is overestimated because the dissociation limit at $R = 7$~bohr
has not fully converged to the separated-atom energy.  This is a grid
convergence issue, not a structural limitation.

% [Placeholder: Figure showing LiH PES comparing no-PK / ab-initio-PK / fitted-PK]


\subsection{BeH$^+$ transferability}
\label{sec:beh}

The same \texttt{ComposedDiatomicSolver} class, with only
$(Z_A, Z_B, M_A, M_B)$ changed, produces a bound BeH$^+$.
Table~\ref{tab:beh} reports the results.

\begin{table}
\caption{BeH$^+$ spectroscopic constants at $l_{\max}=2$.
  Reference values are indicative, from high-level ab initio
  calculations on light hydride ions~\cite{bishop1979}.}
\label{tab:beh}
\begin{ruledtabular}
\begin{tabular}{lddd}
\textrm{Property}
  & \multicolumn{1}{c}{Ab initio PK}
  & \multicolumn{1}{c}{Fitted PK}
  & \multicolumn{1}{c}{Reference} \\
\colrule
$R_{\rm eq}$ (bohr)     & 3.25  & 2.82  & \multicolumn{1}{c}{${\sim}\,2.48$} \\
$\omega_e$ (cm$^{-1}$)  & 2857  & 3088  & \multicolumn{1}{c}{${\sim}\,2222$} \\
Bound?  & \multicolumn{1}{c}{Yes}
  & \multicolumn{1}{c}{Yes}
  & \multicolumn{1}{c}{Yes} \\
\end{tabular}
\end{ruledtabular}
\end{table}

No molecule-specific code or parameter tuning is required.  The PK
parameters are derived from the Be$^{2+}$ core solution using the
identical formulas [Eqs.~(\ref{eq:reff})--(\ref{eq:A})].


\subsection{Sensitivity analysis}
\label{sec:sensitivity}

% [Placeholder: Figure showing R_eq vs A and R_eq vs B]

Scanning the PK parameters independently reveals gradual, monotonic
dependence:
\begin{itemize}
  \item $A$ scan ($B$ fixed at 7.0): $A \in [1, 10]$ maps to an
    $R_{\rm eq}$ range of ${\sim}1.5$~bohr.
  \item $B$ scan ($A$ fixed at 6.93): $B \in [0.5, 10]$ maps to an
    $R_{\rm eq}$ range of ${\sim}1.0$~bohr.
\end{itemize}
The equilibrium is not sharply tuned to specific PK values.
The physics (core screening~+~Pauli repulsion) determines the
equilibrium, not parameter adjustment.


\subsection{$B$ formula comparison}
\label{sec:B_compare}

Four definitions of the core radius were tested, each producing a
different $B = 1/(2r_{\rm core}^2)$ and hence a different LiH
equilibrium.  Table~\ref{tab:B} reports the results.

\begin{table}
\caption{Effect of $B$ formula on LiH $R_{\rm eq}$ at $l_{\max}=2$.
  All use ab initio $A = 6.93$.}
\label{tab:B}
\begin{ruledtabular}
\begin{tabular}{ldddd}
\textrm{Method} & \multicolumn{1}{c}{$r_{\rm core}$ (bohr)}
  & \multicolumn{1}{c}{$B$ (bohr$^{-2}$)}
  & \multicolumn{1}{c}{LiH $R_{\rm eq}$ (bohr)}
  & \multicolumn{1}{c}{Error (\%)} \\
\colrule
RMS                 & 0.677  & 1.09  & 4.00  & 32.7 \\
Peak                & 0.717  & 0.97  & 4.06  & 34.5 \\
Median              & 0.500  & 2.00  & 3.75  & 24.4 \\
Inv-sq (selected)   & 0.267  & 7.00  & 3.21  & 6.4  \\
\end{tabular}
\end{ruledtabular}
\end{table}

The inverse-square formula [Eq.~(\ref{eq:reff})] wins because it
weights the inner core---where the PK barrier must act---via the
$\langle 1/r^2 \rangle$ expectation.  The RMS and peak formulas are
dominated by the diffuse tail and produce barriers that are too broad.
The median is intermediate.  This is not parameter selection---it is
choosing the physically correct moment of the core density for the
purpose at hand (repulsive barrier localization).


%% ========================================================================
%% V. DISCUSSION
%% ========================================================================
\section{Discussion}
\label{sec:discussion}

\subsection{The composed geometry as a category}
\label{sec:category}

The composition $G_{\rm total} = G_{\rm nuc} \ltimes G_{\rm core}
\ltimes G_{\rm val}$ is a fiber bundle with typed inter-layer edges.
This is a mathematical structure (a category) where objects are natural
geometries and morphisms are the continuous deformations
($R \to \infty$ dissociation, $Z \to Z'$ isoelectronic substitution)
that transform one composed geometry into another.  The
LiH$\,{\to}\,$BeH$^+$ transfer is such a morphism: the same functor
(\texttt{ComposedDiatomicSolver}) maps different input parameters to
different molecules.

% [Placeholder: Figure showing the composed graph architecture with
%  three layers: nuclear, core, valence, and typed inter-layer edges]


\subsection{What the $R_{\rm eq}$ improvement means}
\label{sec:req}

The LCAO approach gave $R_{\rm eq} = 2.5$~bohr because the graph
Laplacian's kinetic energy is $R$-independent: the atomic lattice
structure does not change as atoms approach.  The composed geometry
fixes this because $Z_{\rm eff}(r)$ encodes the core's response to the
approaching H atom: at large $R$, the valence sees $Z_{\rm eff} \approx 1$
(screened Li$^+$); at small $R$, the PK barrier prevents core
penetration.  The equilibrium is the balance between bonding
stabilization (from shared valence density) and Pauli repulsion (from
core orthogonalization).  This balance was absent in the LCAO approach.


\subsection{Comparison with standard quantum chemistry}
\label{sec:comparison}

The pseudopotential concept has a long history in solid-state
physics~\cite{heine1970,bachelet1982}.
Standard pseudopotential methods (effective core potentials, ECPs) solve
the same problem within Gaussian basis set frameworks.  Our approach is
structurally different: the core is solved in its natural geometry
(hyperspherical), the valence in its natural geometry (molecule-frame
hyperspherical), and the coupling is through the core wavefunction
itself rather than through fitted ECP parameters.  The analogy:
ECPs are to our ab initio PK as Gaussian basis sets are to our
natural geometry lattices---both solve the same physics from
different starting points.


\subsection{Limitations}
\label{sec:limitations}

\begin{enumerate}
\item The Gaussian$/r^2$ PK form [Eq.~(\ref{eq:Vpk})] is approximate.
  The exact Phillips--Kleinman pseudopotential is nonlocal and
  energy-dependent~\cite{pk1959}.  Our Gaussian approximation captures
  the dominant repulsion but not its angular or energy dependence.

\item Core--valence exchange integrals are omitted.  These would
  contribute ${\sim}1$--5\% corrections to the binding energy.

\item Only $\sigma$ channels at $l_{\max} = 2$.  Paper~15 showed that
  $\pi$ channels add ${\sim}7$~percentage points for H$_2$; similar
  improvement is expected for LiH at $l_{\max} \geq 3$.

\item The dissociation limit is not converged at $R = 7$~bohr.
  $D_e$ values are upper bounds.

\item Only two-core plus two-valence systems.  Extension to more
  electrons requires composing additional Level~3/4 subproblems (e.g.,
  lone pairs as Level~3, additional bonds as Level~4).
\end{enumerate}


\subsection{Outlook: the composition pattern for larger molecules}
\label{sec:outlook}

The composed geometry generalizes:
\begin{itemize}
  \item \textbf{BeH$_2$:} 1~core (Be~$1s^2$) + 2~bond pairs (Level~4
    on each Be--H) + inter-bond coupling.
  \item \textbf{H$_2$O:} 1~core (O~$1s^2$) + 2~bond pairs (Level~4 on
    each O--H) + 2~lone pairs (Level~3 on O) + inter-pair coupling.
  \item \textbf{General:} partition electrons into cores (Level~3 per
    atom) + bond pairs (Level~4 per bond) + lone pairs (Level~3 on atoms).
    The molecule is a \emph{hypergraph} where each hypernode is a
    natural-geometry subproblem and each hyperedge is a typed coupling.
\end{itemize}

The immediate next step is LiH at $l_{\max} = 4$ with $\sigma{+}\pi$
channels, which should reduce the $R_{\rm eq}$ error from 6.4\% to
${\sim}2$--3\% based on the $l_{\max}$ convergence pattern in
Paper~15.


%% ========================================================================
%% VI. CONCLUSION
%% ========================================================================
\section{Conclusion}
\label{sec:conclusion}

Composed natural geometries solve the core-valence problem without a
single high-dimensional coordinate system.  The key results:

\begin{enumerate}
\item LiH $R_{\rm eq}$ improves from 17\% error (LCAO) to 6.4\%
  (composed, ab initio PK) with zero molecular fitting.

\item The ab initio PK pseudopotential is derived entirely from the
  core wavefunction: $A$ from the core--valence energy gap
  [Eq.~(\ref{eq:A})], $B$ from the $\langle 1/r^2 \rangle$-weighted
  core radius [Eq.~(\ref{eq:reff})].

\item The architecture transfers to BeH$^+$ with no per-molecule
  tuning.

\item The composition pattern---core graphs + valence graphs + typed
  coupling edges---provides a concrete framework for extending discrete
  graph quantum chemistry to general molecules.
\end{enumerate}

The natural geometry hierarchy is now:
\begin{center}
\begin{tabular}{cccc}
\toprule
Level & System & Geometry & Best result \\
\midrule
1 & H       & $S^3$ (Fock)           & $<0.1\%$ error \\
2 & H$_2^+$ & Prolate spheroid       & 0.70\% error \\
3 & He      & Hyperspherical         & 0.05\% error \\
4 & H$_2$   & Mol-frame hypersp.     & 94.1\% $D_e$ \\
5 & LiH     & Composed (Level 3+4)   & $R_{\rm eq}$: 6.4\% error \\
5 & BeH$^+$ & Composed (Level 3+4)   & Bound, physical \\
\bottomrule
\end{tabular}
\end{center}


%% ========================================================================
%% REFERENCES
%% ========================================================================
\begin{thebibliography}{25}

\bibitem{fock1935}
V.~A.~Fock,
``Zur Theorie des Wasserstoffatoms,''
Z.~Phys.\ \textbf{98}, 145 (1935).

\bibitem{pk1959}
J.~C.~Phillips and L.~Kleinman,
``New method for calculating wave functions in crystals and molecules,''
Phys.\ Rev.\ \textbf{116}, 287 (1959).

\bibitem{lin1995}
C.~D.~Lin,
``Hyperspherical coordinate approach to atomic and other
Coulombic three-body systems,''
Phys.\ Rep.\ \textbf{257}, 1 (1995).

\bibitem{huber1979}
K.~P.~Huber and G.~Herzberg,
\emph{Molecular Spectra and Molecular Structure.\ IV.\ Constants of
Diatomic Molecules}
(Van Nostrand Reinhold, New York, 1979).

\bibitem{nist}
A.~Kramida, Yu.~Ralchenko, J.~Reader, and NIST ASD Team,
NIST Atomic Spectra Database (ver.~5.11),
\url{https://physics.nist.gov/asd} (2024).

\bibitem{heine1970}
V.~Heine and D.~Weaire,
``Pseudopotential theory of cohesion and structure,''
Solid State Phys.\ \textbf{24}, 249 (1970).

\bibitem{bachelet1982}
G.~B.~Bachelet, D.~R.~Hamann, and M.~Schl\"uter,
``Pseudopotentials that work: From H to Pu,''
Phys.\ Rev.\ B \textbf{26}, 4199 (1982).

\bibitem{bishop1979}
D.~M.~Bishop and L.~M.~Cheung,
``A theoretical investigation of HeH$^+$,''
J.~Mol.\ Spectrosc.\ \textbf{75}, 462 (1979).

\bibitem{paper0}
J.~Loutey,
``Geometric Packing and the Universal Constant''
(Paper~0 in the Geometric Vacuum series).

\bibitem{paper1}
J.~Loutey,
``Spectral Graph Theory for Atomic Structure''
(Paper~1 in the Geometric Vacuum series).

\bibitem{paper7}
J.~Loutey,
``The Dimensionless Vacuum: Scale-Invariant Graph Topology
and the Schr\"odinger Equation''
(Paper~7 in the Geometric Vacuum series).

\bibitem{paper10}
J.~Loutey,
``Nuclear Lattice: Rovibrational Spectra from SU(2) Algebraic
Chains'' (Paper~10 in the Geometric Vacuum series).

\bibitem{paper11}
J.~Loutey,
``Molecular Fock Projection: Prolate Spheroidal Lattice
for Diatomic Molecules''
(Paper~11 in the Geometric Vacuum series).

\bibitem{paper12}
J.~Loutey,
``Algebraic Two-Electron Integrals: Neumann Expansion
for $V_{ee}$ on the Prolate Spheroidal Lattice''
(Paper~12 in the Geometric Vacuum series).

\bibitem{paper13}
J.~Loutey,
``Hyperspherical Lattice: Two-Electron Atoms as
Coupled-Channel Graphs''
(Paper~13 in the Geometric Vacuum series).

\bibitem{paper15}
J.~Loutey,
``The Level~4 Natural Geometry: Two-Center Two-Electron Molecules
as Hyperspherical Fiber Bundles''
(Paper~15 in the Geometric Vacuum series).

\bibitem{paper16}
J.~Loutey,
``Chemical Periodicity from $S_N$ Representation Theory''
(Paper~16 in the Geometric Vacuum series).

\bibitem{fci_a}
J.~Loutey,
``Full Configuration Interaction on the Geometric Lattice: Atoms''
(FCI-A paper in the Geometric Vacuum series).

\bibitem{fci_m}
J.~Loutey,
``Full Configuration Interaction on the Geometric Lattice: Molecules''
(FCI-M paper in the Geometric Vacuum series).

\bibitem{geovac2026}
J.~Loutey,
``GeoVac: Computational Quantum Chemistry via Spectral Graph Theory,''
version 1.6.1 (2026),
\url{https://github.com/jloutey-hash/geovac}.

\end{thebibliography}

\end{document}
